\documentclass[11pt]{article}
\usepackage[margin=1in]{geometry}
\usepackage{amsmath}
\usepackage{booktabs}
\usepackage[hidelinks]{hyperref}

\newcommand{\TyphoonPending}[1]{\textbf{[PENDING: #1]}}

\title{Event-Resolved Gravimetric Detection of Typhoon-Driven Ocean Mass Redistribution}
\author{\TyphoonPending{author list and data-provider collaboration}}
\date{Methods-first draft; not for submission}

\begin{document}
\maketitle

\begin{abstract}
We test whether named typhoons produce event-resolved ocean mass-redistribution
signals in superconducting-gravimeter records after explicit separation of
atmosphere, hydrology, direct ocean attraction, elastic loading, and instrument
effects. The design combines gravimetry with pressure, best tracks, gridded sea
level, tide gauges and GNSS, and evaluates a completely held-out typhoon together
with strong-weather and quiet controls. \TyphoonPending{Insert data-gated
multi-event results and uncertainty after station authorization.}
\end{abstract}

\section{Introduction}
Previous work has modeled non-tidal ocean loading in the South China Sea and
observed storm-surge loading at Helgoland. The increment tested here is not
general NTOL correction, but named-typhoon attribution with a closed component
budget, independent evidence, and cross-event generalization.

\section{Data and preregistered event design}
The raw-data gate requires either one station covering at least three typhoons or
two stations covering at least two typhoons each, including gravity, colocated
pressure, calibration/timing/jump/maintenance records, sea-level fields, storm
tracks, precipitation, and hydrology. Evaluation additionally requires one
held-out typhoon, one non-typhoon strong-weather control, three quiet windows,
and no overlapping same-station windows across splits.

IBTrACS agency wind and pressure fields are retained separately. ERA5 surface
pressure supplies atmospheric loading; mean-sea-level pressure and winds are
event diagnostics. Significant wave height is not converted into additional
water mass. CMEMS sea level is admitted only after its tide, mean dynamic
topography, and inverse-barometer semantics are audited.

\section{Methods}
\subsection{Transparent gravimeter correction}
For observed gravity $g_{\mathrm{obs}}$, the residual is
\begin{equation}
 g_{\mathrm{res}} = g_{\mathrm{obs}}-
 \sum_k g_k,
\end{equation}
where every removed component has a unique physical-effect identifier, source,
SI series, and pre-applied flag. Two components cannot own the same effect. Each
stage stores its input, removed series, output, and peak contribution.

Missing samples and cadence violations split records without interpolation.
Candidate discontinuities are flagged at original indices but are not
automatically classified or corrected.

\subsection{Independent ocean forward model}
Sea-level anomaly $\eta$ is converted to surface density
$\Delta\sigma=\rho_w\eta$. Direct attraction is integrated on spherical and
WGS~84 grids with explicit masks and partial coastal-cell fractions. The modeled
ocean signal is
\begin{equation}
 g_{\mathrm{ocean}}^{\mathrm{model}}=
 g_{\mathrm{direct}}+g_{\mathrm{deformation}}+g_{\mathrm{internal}},
\end{equation}
with every term retained separately. Elastic terms are disabled until the
Green-function Earth model, reference frame, normalization, and component
semantics pass a published benchmark.

\subsection{Double-count audit}
A CMEMS field containing inverse-barometer response cannot be combined with a
second ERA5-driven inverse-barometer correction. Likewise, a pre-corrected NTOL
residual cannot receive NTOL again. Unknown product physics blocks assembly.

\subsection{Evaluation}
Primary metrics are correlation, RMSE, peak amplitude and time error, explained
variance, frequency-domain coherence, event SNR, and held-out performance.
Quiet/control windows estimate false positives. Model, filter, hydrology,
Green-function, station, and event sensitivities are retained, including failed
events.

\section{Results}
A registered non-restricted readiness audit connects 17 completed processing
and evaluation work units. The open IBTrACS inventory contains 48 broad regional
track candidates and 11 priority station-coverage inquiries. These are not
gravity events and do not select a pilot: the real-event gate remains
\texttt{pending\_no\_event\_windows}. The environmental-effect ledger also
remains open because the exact historical CMEMS inverse-barometer semantics are
unknown. Consequently the claim-safe manuscript branch is
\texttt{pending\_evidence}. \TyphoonPending{Insert only registered real-SG
event results after the station and effect-ownership gates pass.}

\section{Discussion}
Possible outcomes include successful event attribution, evidence that ocean
products miss rapid coastal redistribution, or an upper constraint from
non-detection. A single successful event cannot establish generality.

\section{Data and code availability}
Restricted SG and station logs will remain outside Git. Versioned manifests will
record owner, terms, station, channels, time range, product version, retrieval
time, checksum, exclusions, and redistribution limits. Processing code and
compact nonrestricted metrics will be released.

\section{Limitations}
SG event windows, calibration and station logs remain authorization-dependent.
The exact historical CMEMS pressure-response semantics are a separate product-
metadata blocker. Consequently no pilot event, attribution coefficient,
held-out improvement, event SNR or false-positive rate is reported. The current
method audit cannot substitute for those observations.

\bibliographystyle{plain}
\bibliography{references}
\end{document}
